\section{Conclusion}

This paper presented an energy-aware offline multimodal SNN pipeline for water-quality prediction in an Amazonian edge-monitoring context, demonstrating practical feasibility from training to deployment, including geospatial interpretation and cross-platform latency/energy benchmarking. Neuromorphic inference proved deployable under real hardware constraints, with platform-specific throughput--energy trade-offs; while not always best in absolute MAE, the SNN remains a robust option when event-driven processing and edge compatibility are first-class requirements.

As discussed above, validity remains bounded by the limited field dataset --- mitigated by repeated-seed evaluation, conservative augmentation, and per-target reporting --- and by hydrological diversity across Amazonian subregions. pH showed the most reliable predictive behavior, turbidity required wider interpretation margins because of its larger absolute error and dispersion, and OD remained the main explained-variance limitation due to negative mean $R^2$ across all evaluated models. Future work includes broader acquisition across seasonal regimes, external validation in additional freshwater systems and lakes, and deeper co-optimization of encoding parameters and hardware execution modes. Planned extensions include full k-fold cross-validation on the enlarged dataset, comparison of augmentation strategies beyond Gaussian perturbation, energy measurement of SNN, Random Forest, LSTM, and MLP under the same idle/load protocol on identical devices, and enriched input modalities to address the dissolved-oxygen limitation. Overall, neuromorphic models enable energy-aware edge intelligence for practical environmental monitoring in remote Amazonian regions.
